\section{Universal and existential quantifiers}\label{sec:03-propositions-and-proofs:06-quantifiers:1}

\begin{lstlisting}
theorem all_nats_ge_zero : ∀ n : Nat, n ≥ 0 :=
  fun n => Nat.zero_le n
\end{lstlisting}

\begin{itemize}
\tightlist
\item
  \texttt{∀ x : α, P x} is again just a (dependent) function type: given any \texttt{x}, produce a proof of \texttt{P x}. \texttt{all\textbackslash{}\_nats\textbackslash{}\_ge\textbackslash{}\_zero} is literally a function taking any \texttt{n} and returning a proof that \texttt{n ≥ 0} --- here, \texttt{Nat.zero\textbackslash{}\_le n} already proves exactly that for whichever \texttt{n} gets passed in.
\end{itemize}

\begin{lstlisting}
theorem exists_even : ∃ n : Nat, n % 2 = 0 :=
  ⟨2, rfl⟩
\end{lstlisting}

\begin{itemize}
\tightlist
\item
  \texttt{∃ x : α, P x} is a structure: a \textbf{witness} value plus a proof that the witness satisfies \texttt{P}.
\end{itemize}

\textbf{Reading the shortcut \texttt{⟨2, rfl⟩}.} This is the same ``here are the pieces, in order'' anonymous-constructor shorthand from

It was already reused for \texttt{∧} (\hyperref[sec:03-propositions-and-proofs:05-and-or-not:1]{Chapter 3 §5}'s \texttt{⟨hp, hq⟩}), and here it builds an \texttt{∃}-proof in one line instead of two. The pattern to remember: \textbf{first the number, then why it works.}

\begin{itemize}
\tightlist
\item
  \texttt{2} is the \textbf{witness} --- the specific number being claimed even. It is the smallest one that actually makes the point (\texttt{0} technically works too, but reads as a trick --- ``of course \texttt{0} is even'').
\item
  \href{https://lean-lang.org/doc/reference/latest/Tactic-Proofs/Tactic-Reference/}{\texttt{rfl}} is the \textbf{proof}, filling in for \texttt{P 2}, i.e.~for \texttt{2 \textbackslash{}\% 2 = 0}. It works by \texttt{rfl} alone (no lemma needed) because \texttt{2 \textbackslash{}\% 2} simply \emph{computes} to \texttt{0} --- both sides of the equation are already the same term once evaluated, exactly like \texttt{2 + 2 = 4} back in \hyperref[sec:03-propositions-and-proofs:01-prop:1]{Chapter 3 §1}.
\end{itemize}

So \texttt{⟨2, rfl⟩ : ∃ n : Nat, n \textbackslash{}\% 2 = 0} reads as ``2 works, and here is why: it just computes.'' Wherever \texttt{⟨w, p⟩} proves an \texttt{∃}-statement elsewhere in this book, it should be read the same way: \texttt{w} is \emph{what} is claimed to exist, and \texttt{p} is \emph{why} it satisfies the property. They are packed together because that is exactly the two pieces of data an existence proof needs.

\begin{lstlisting}
-- A minimal, self-contained primality check (no Mathlib import needed):
-- `n` is prime if it's at least 2 and no number strictly between 2 and n
-- divides it. `@[reducible]` lets `decide` below see straight through
-- this definition, instead of needing a separate `unfold` step first.
@[reducible] def isPrime (n : Nat) : Prop :=
  n ≥ 2 ∧ ∀ m : Nat, m < n → m ≥ 2 → ¬ (m ∣ n)

theorem exists_prime_gt_three : ∃ p : Nat, p > 3 ∧ isPrime p :=
  ⟨5, by decide⟩
\end{lstlisting}

\textbf{A second example: there is always a bigger prime.} \texttt{exists\textbackslash{}\_prime\textbackslash{}\_gt\textbackslash{}\_three} is one concrete instance of a classical fact: no matter which number is chosen, there is a prime bigger than it (here the chosen number is \texttt{3}, and \texttt{5} is a prime that beats it). It follows the exact same ``number, then why'' shape as before, but with two differences worth noting:

\begin{itemize}
\tightlist
\item
  The \emph{property} being witnessed, \texttt{p > 3 ∧ isPrime p}, is itself an \texttt{∧}. So the full picture is a witness (\texttt{5}) plus a \emph{pair} of facts about it (\texttt{5 > 3}, and \texttt{5} is prime), all packed into the outer \texttt{⟨\textbackslash{}\_, \textbackslash{}\_⟩}.
\item
  The proof is not \texttt{rfl}. \texttt{p > 3 ∧ isPrime p} does not reduce to a plain equality, so instead the second slot is by \href{https://lean-lang.org/doc/reference/latest/Tactic-Proofs/Tactic-Reference/}{\texttt{decide}} --- the same brute-force tactic from \hyperref[sec:03-propositions-and-proofs:05-and-or-not:1]{Chapter 3 §5}'s \texttt{not\textbackslash{}\_example}, which here checks \texttt{5 > 3} outright and tries every candidate divisor below \texttt{5} to confirm none of them divide it.
\end{itemize}

Read \texttt{⟨5, by decide⟩} exactly as before: \texttt{5} is still just the witness, \texttt{by decide} is still just the proof, filled in by a tactic block instead of \texttt{rfl} because \emph{this} proof happens to need one. The witness-then-proof shape never changes; only \emph{how} the proof half gets produced does.

This one example only shows \emph{a} prime past \texttt{3}. It does not yet show that this works for \emph{every} number that could have been chosen instead of \texttt{3}. That stronger claim, \(\forall n,\ \exists p > n,\ \mathrm{isPrime}\ p\) (Euclid's theorem, first proved around 300 BCE), needs an argument that works uniformly for an \emph{arbitrary} \texttt{n}. That means induction, not yet introduced. \hyperref[sec:04-tactics:00-index:1]{Chapter 4} covers the tactics that make an argument like that possible. A fully formalized proof of Euclid's theorem lives in Mathlib as \texttt{Nat.exists\textbackslash{}\_infinite\textbackslash{}\_primes}, one of the ``Read more'' pointers once \hyperref[sec:13-next-steps:00-index:1]{Chapter 13} introduces Mathlib itself.

\textbf{Remark (a more formal restatement).} The bullet points above already state everything needed to use \texttt{∃} and \texttt{∀}. This paragraph and the ``Mathematical reading'' box below only restate the same facts in more formal language, for readers who want the connection made fully explicit. \texttt{∃ x : α, P x} being ``a witness plus a proof'' is, formally, saying it is a \textbf{dependent pair type}: the type of the second component (the proof) depends on the \emph{value} chosen for the first (the witness). It would be \texttt{P 2} for witness \texttt{2}, but \texttt{P 1} for witness \texttt{1}, a genuinely different type each time. This dependency is exactly why \texttt{∃} cannot be built from the ordinary (non-dependent) pairing used for \texttt{∧}, and why Lean needs a dedicated \texttt{Exists} former for it.

\begin{mathreading}

The two quantifiers are the ``indexed'' versions of the product and sum seen for \(\wedge\) and \(\vee\) in the previous section. Instead of combining two fixed propositions \(P\) and \(Q\), they combine a whole \emph{family} of propositions \(P(x)\), one for each \(x \in \alpha\). Universal quantification is the (\(\Pi\)-)type \[
\forall x{:}\alpha,\ P(x) \;=\; \prod_{x : \alpha} P(x),
\] literally an \(\alpha\)-indexed product: a proof is a function assigning to each \(x\) a proof of \(P(x)\), so \texttt{all\textbackslash{}\_nats\textbackslash{}\_ge\textbackslash{}\_zero} is the map \(n \mapsto (0
\le n)\). This exactly generalizes how a proof of \(P \wedge Q\) was a pair \((p, q)\), one component per conjunct, except now there is one component per element of \(\alpha\) rather than just two. Existential quantification is the (\(\Sigma\)-)type \[
\exists x{:}\alpha,\ P(x) \;=\; \sum_{x : \alpha} P(x),
\] an \(\alpha\)-indexed sum: a proof is a dependent pair \(\langle a, h\rangle\) with \(a \in \alpha\) the witness and \(h : P(a)\). Here \(\langle 2,
\mathrm{refl}\rangle\) witnesses \(2 \bmod 2 = 0\). This is the \emph{constructive} reading of \(\exists\): asserting existence requires exhibiting an explicit witness \(a\) together with a proof \(h\) that it works. This exactly generalizes how a proof of \(P \vee Q\) was a tagged choice, except now the ``tag'' is which element of \(\alpha\) was chosen.

\end{mathreading}
